% ============================================================================
% SECTION 5: CPA AS BAYESIAN PROJECTION DYNAMICS
% ============================================================================
\section{Correlation Power Analysis as Bayesian Projection Dynamics}\label{sec:cpa}

\subsection{Leakage model}\label{subsec:adv-model}

We consider the standard non-profiled CPA setting.  The adversary knows $S$,
observes plaintexts $a_1,\ldots,a_n$ uniformly drawn from $\B{s}$, and obtains
leakage measurements
\begin{equation}\label{eq:leakage-model}
  W_i = f(S(a_i\oplus k^*)) + \varepsilon_i,
  \qquad \varepsilon_i \sim \mathcal N(0,\sigma^2),
\end{equation}
where $k^*$ is the true key and $f:\B{s}\to\R$ is the leakage model.  In the
numerical experiments we take $f = \HW$.

\subsection{Bayesian update as a mixture projection}\label{subsec:bayes-m}

\begin{theorem}[CPA posterior update as $m$-projection]\label{thm:cpa-m-proj}
  Let $q_{i-1}\in\Delta_{2^s-1}^{\circ}$ be the prior belief over keys before
  observing $(a_i,W_i)$.  The posterior
  \begin{equation}\label{eq:bayes-update}
    q_i(k) = \frac{q_{i-1}(k)\,p(W_i\mid a_i,k)}
                 {\sum_{h\in\B{s}} q_{i-1}(h)\,p(W_i\mid a_i,h)}
  \end{equation}
  is the $m$-projection of the joint prior-likelihood distribution onto the
  affine constraint defined by the observation, marginalised to the key
  coordinate.
\end{theorem}

\begin{proof}
  The joint distribution before conditioning is $q_{i-1}(k)p(a_i)p(W_i\mid a_i,k)$.
  Conditioning on the observed value imposes an affine constraint.  Minimising
  the Kullback‑Leibler divergence under that constraint yields the conditional
  distribution, whose key marginal is exactly~\eqref{eq:bayes-update}.
\end{proof}

The exact S‑box distributions $P_k$ lie on the boundary of the simplex, but
the belief trajectory $q_0,q_1,\ldots,q_n$ stays in the interior of the belief
simplex as long as the prior has full support and the Gaussian likelihood is
positive.  Hence CPA can be studied as a regular trajectory on $(Q,g_F^{(Q)})$.

\subsection{Per-trace signal information}\label{subsec:per-trace}

\begin{definition}[Per-trace signal information]\label{def:per-trace}
  For a leakage model $f$ and noise variance $\sigma^2$, define
  \begin{equation}\label{eq:per-trace-def}
    I_S(f,\sigma) = \sigma^{-2}
    \operatorname{Var}_{a\sim\mathrm{Unif}(\B{s})}\bigl[f(S(a\oplus k^*))\bigr].
  \end{equation}
  When $S$ is bijective, this quantity does not depend on $k^*$.
\end{definition}

\begin{proposition}[Hamming-weight specialization]\label{thm:per-trace}
  For $f = \HW$ and a bijective $S$,
  \begin{equation}\label{eq:hw-info}
    I_S(\HW,\sigma) = \frac{s}{4\sigma^2}.
  \end{equation}
  In particular, for 4‑bit bijective S‑boxes, $I_S(\HW,\sigma)=1/\sigma^2$.
\end{proposition}

\begin{proof}
  Bijectivity of $S$ together with the uniformity of $a\oplus k^*$ implies that
  $S(a\oplus k^*)$ is uniform over $\B{s}$.  The Hamming weight of a uniform
  $s$-bit word follows a binomial distribution $\mathrm{Binomial}(s,1/2)$, whose
  variance is $s/4$.
\end{proof}

\subsection{Pairwise log-likelihood increments}\label{subsec:pairwise-llr}

For a competing key $k \neq k^*$, the per‑trace squared separation is
\begin{equation}\label{eq:delta-k}
  \Delta(k,k^*) = \frac{1}{2^s} \sum_{a\in\B{s}}
  \bigl(f(S(a\oplus k^*)) - f(S(a\oplus k))\bigr)^2.
\end{equation}

\begin{proposition}[Expected pairwise log‑Bayes gain]\label{prop:pairwise-gain}
  Under the Gaussian leakage model,
  \begin{equation}\label{eq:llr-gain}
    \mathbb E\left[\log\frac{p(W\mid A,k^*)}{p(W\mid A,k)}\right]
    = \frac{\Delta(k,k^*)}{2\sigma^2},
  \end{equation}
  where the expectation is over $A$ (uniform) and $W$ conditioned on $A$ and $k^*$.
\end{proposition}

\begin{proof}
  For a fixed plaintext $a$, the two likelihoods are Gaussian with the same
  variance $\sigma^2$ and means $\mu_*(a)=f(S(a\oplus k^*))$,
  $\mu_k(a)=f(S(a\oplus k))$.  The Kullback‑Leibler divergence between two
  univariate Gaussians with equal variance is $(\mu_*(a)-\mu_k(a))^2/(2\sigma^2)$.
  Averaging over $a$ gives the result.
\end{proof}

The worst‑case pairwise gap
\begin{equation}\label{eq:gamma-gap}
  \Gamma_S(f,\sigma) = \min_{k\neq k^*} \frac{\Delta(k,k^*)}{2\sigma^2}
\end{equation}
provides a rigorous likelihood‑separation measure, distinct from $I_S(f,\sigma)$
in general.

\subsection{Model-based sample-complexity estimate}\label{subsec:convergence}

Starting from a uniform prior, the initial key uncertainty equals $s\log2$ nats.
A first‑order information estimate then yields
\begin{equation}\label{eq:sample-estimate}
  n_{\mathrm{est}} \approx \frac{s\log2}{I_S(f,\sigma)}.
\end{equation}
For 4‑bit Hamming‑weight leakage, $n_{\mathrm{est}}\approx 4\log2\,\sigma^2$.
This is a heuristic leakage‑rate estimate.  A more rigorous bound follows from
the pairwise gap $\Gamma_S(f,\sigma)$: the expected log‑posterior odds against
any fixed wrong key grow at rate at least $\Gamma_S(f,\sigma)$.
